\section{指标设置}

为全面评价 AGF-ADNet 在多采样率异常检测任务中的判别性能，本文将异常检测问题视为二分类任务，并从总体判别能力、异常识别能力、误报控制能力以及综合性能四个方面设置评价指标。具体而言，本文采用准确率、精确率、召回率、故障检测率、误报率和 F1-Score 六项指标。上述指标能够从不同角度刻画模型在正常样本与异常样本上的判别表现，避免单一指标导致的评价片面性。

\subsection{混淆矩阵定义}

设测试样本的真实标签为 $y$，模型预测标签为 $\hat{y}$。在二分类异常检测任务中，通常将异常状态记为正类，将正常状态记为负类。由此可定义四种基本判别结果：
\begin{itemize}
    \item 真正例（True Positive, TP）：真实为异常且模型预测为异常的样本数；
    \item 真负例（True Negative, TN）：真实为正常且模型预测为正常的样本数；
    \item 假正例（False Positive, FP）：真实为正常但模型预测为异常的样本数；
    \item 假负例（False Negative, FN）：真实为异常但模型预测为正常的样本数。
\end{itemize}

基于上述四个统计量，可以构造本文所采用的六项评价指标。由于工业异常检测通常兼具类别不均衡、异常代价高和误报敏感等特点，因此有必要同时考察对异常样本的检出能力与对正常样本的误判情况。

\subsection{准确率}

准确率（Accuracy）用于衡量模型对全部样本的总体判别正确程度，其定义为
\begin{equation}
    \mathrm{Accuracy}
    =
    \frac{TP + TN}{TP + TN + FP + FN}.
\end{equation}

该指标反映模型在整体层面上的分类正确比例，数值越大表示总体识别性能越好。准确率能够直接体现模型对测试集的全局判别能力，因此是最基础的评价指标之一。

然而，在异常检测任务中，若正常样本显著多于异常样本，则准确率可能受到类别分布的影响。例如，当异常样本占比较低时，即使模型对异常状态识别不足，只要能够正确判别大多数正常样本，仍可能获得较高的准确率。因此，准确率适合作为总体性能指标，但不能单独用于判断模型对异常状态的识别效果。

\subsection{精确率}

精确率（Precision）用于衡量模型判定为异常的样本中真正异常样本所占的比例，其定义为
\begin{equation}
    \mathrm{Precision}
    =
    \frac{TP}{TP + FP}.
\end{equation}

精确率越高，说明模型输出的异常报警中真实异常所占比例越高，也即误报程度越低。在工业监测场景中，异常报警往往会触发进一步巡检、停机检查或调度干预，因此较高的精确率意味着模型所发出的报警更具有可信度，能够减少不必要的人工排查和生产扰动。

需要指出的是，精确率强调的是报警结果的可靠性，而不是异常样本是否被尽可能完整地检出。因此，当模型采取较为保守的报警策略时，精确率可能较高，但同时也可能遗漏部分真实异常。基于这一原因，精确率通常需要与召回率等指标联合分析。

\subsection{召回率}

召回率（Recall）用于衡量全部真实异常样本中被模型成功识别出来的比例，其定义为
\begin{equation}
    \mathrm{Recall}
    =
    \frac{TP}{TP + FN}.
\end{equation}

召回率越高，表示模型对异常状态的覆盖能力越强，漏检的异常样本越少。对于电源系统和工业过程监测任务而言，漏检往往意味着未能及时发现潜在故障，从而可能造成控制性能恶化、设备损伤甚至系统失稳。因此，召回率是异常检测中极为重要的核心指标。

与精确率相比，召回率更关注对真实异常的完整检出能力。当模型倾向于扩大报警范围时，召回率通常会上升，但这也可能伴随更多误报。因此，召回率虽能体现模型对异常样本的敏感性，却不能独立反映报警质量。

\subsection{故障检测率}

故障检测率（Fault Detection Rate, FDR）用于衡量模型对故障或异常状态的检出能力，其定义为
\begin{equation}
    \mathrm{FDR}
    =
    \frac{TP}{TP + FN}.
\end{equation}

在二分类异常检测设定中，若将异常样本视为正类，则故障检测率与召回率在数值形式上完全一致。本文同时保留召回率与故障检测率两个指标，主要是出于学术表述和工程评价习惯的双重考虑。召回率更常用于统计学习和模式识别领域，用以描述正类样本的检出比例；故障检测率则更符合故障诊断与工业过程监测领域的术语体系，用以直接反映异常工况的发现能力。

因此，在本文实验中，召回率和故障检测率的计算结果相同，但二者的解释侧重点有所不同。前者突出模型对异常正类的覆盖程度，后者强调模型在工程场景中的故障发现能力。

\subsection{误报率}

误报率（False Alarm Rate, FRA）用于衡量真实正常样本中被错误判定为异常的比例，其定义为
\begin{equation}
    \mathrm{FRA}
    =
    \frac{FP}{FP + TN}.
\end{equation}

误报率越低，说明模型越不容易将正常波动、采样不均匀或轻微噪声误判为异常。在工业异常检测中，误报过多会显著降低监测系统的可用性和操作人员对报警结果的信任度，甚至可能造成频繁且无效的人工干预。因此，误报率是评估模型工程可部署性的重要指标。

与精确率不同，误报率的分母是全部真实正常样本，因此它更直接反映模型对正常工况的保持能力。若一个模型在异常检出方面表现较好，但误报率过高，则说明其报警边界过于宽松，在实际应用中仍然存在较大局限。

\subsection{F1-Score}

F1-Score 用于综合衡量精确率与召回率之间的平衡关系，其定义为
\begin{equation}
    \mathrm{F1}
    =
    \frac{2 \times \mathrm{Precision} \times \mathrm{Recall}}
    {\mathrm{Precision} + \mathrm{Recall}}
    =
    \frac{2TP}{2TP + FP + FN}.
\end{equation}

F1-Score 本质上是精确率和召回率的调和平均，因此只有当二者都保持较高水平时，F1-Score 才会取得较大数值。若模型虽然能够检出较多异常但伴随大量误报，则精确率下降会导致 F1-Score 降低；若模型报警非常谨慎而漏检较多异常，则召回率下降同样会拉低 F1-Score。

对于多采样率异常检测任务而言，F1-Score 能够较好地反映模型在异常发现能力与报警可靠性之间的综合平衡状态，因此常作为最终性能比较中的关键指标之一。

\subsection{六项指标的综合作用}

上述六项指标各自关注的性能侧重点并不相同。准确率反映总体判别正确程度，精确率反映异常报警的可信度，召回率与故障检测率反映异常样本的检出能力，误报率反映对正常样本的保持能力，F1-Score 则综合评价精确率与召回率之间的均衡关系。

在本文的实验分析中，六项指标联合使用具有以下意义：
\begin{enumerate}
    \item 可从总体性能与异常识别性能两个层面同时评价模型；
    \item 可同时考察漏检与误报两类风险，避免只关注单一错误类型；
    \item 可兼顾统计学习领域与故障诊断领域常用术语体系，增强结果表述的完整性；
    \item 可为不同模型之间的横向比较提供更加全面的依据。
\end{enumerate}

因此，本文选取准确率、精确率、召回率、故障检测率、误报率和 F1-Score 作为实验部分的主要评价指标，以系统刻画 AGF-ADNet 在多采样率异常检测任务中的综合性能。
